\subsection*{Del libro a un conjuro listo: ejemplo completo}
\begin{enumerate}
\item Entra en CAMP y elige Memorizar Conjuros para un mago, u Orar Conjuros para un clérigo o paladín.
\item Selecciona el personaje y el nivel de conjuro. Añade copias con clic izquierdo o el botón +; quítalas con clic derecho o el botón $-$. Si indica 0 restantes, primero libera un hueco. Por ejemplo, cambia una Ayuda por un Paralizar.
\item Sal de la lista y elige Descansar grupo. Elegir conjuros no los recupera al instante. Espera a que los mensajes indiquen que han sido preparados; si interrumpes el descanso puede quedar trabajo pendiente.
\item Vuelve a la aventura, usa el libro o símbolo desde una mano, elige el nivel y el conjuro listo. Si pide un destinatario, señala su retrato.
\item Tras gastarlo, vuelve a descansar para recuperarlo. Para aprender una magia nueva de mago, primero copia el pergamino y después selecciónala para preparar.
\end{enumerate}

\textbf{Teclado en la preparación.} Las flechas arriba y abajo resaltan un conjuro; izquierda y derecha cambian de nivel. Enter o espacio añade una copia, la tecla menos quita una y Esc sale conservando la selección. El número de la derecha es la cantidad elegida, no un campo donde escribir cifras.

\textbf{Contacto, proyectil y protección.} Un proyectil necesita una dirección libre hacia el objetivo; un conjuro de contacto puede exigir primera fila o una mano libre. En Golpe eléctrico, Toque vampírico y Hoja llameante debes usar después la mano que contiene la magia. Una protección no hace daño por sí sola. La sección 13 explica cada caso.

\textbf{Ritmo del combate.} Ataca, espera a que la mano vuelva a estar disponible y vuelve a usarla. Retroceder o desplazarte de lado puede darte margen, siempre que hayas comprobado el suelo. No abras el inventario esperando detener a los enemigos. Recupera los proyectiles que hayan quedado en el suelo cuando la zona sea segura.
